% 1. Preamble
\documentclass[12pt, a4paper]{article}

% Import Packages
\usepackage[utf8]{inputenc}        % Encoding
\usepackage{amsmath, amssymb}     % Advanced math formulas
\usepackage{mathtools}           % Math enhancements
\usepackage{graphicx}            % Image insertion
\usepackage{hyperref}            % Interactive hyperlinks & PDF bookmarks
\usepackage{color}               % Color text and backgrounds
\usepackage{geometry}            % Page geometry
\usepackage{comment}             % Multi-line comments
\usepackage{titling}             % Custom title formatting

\geometry{a4paper, margin=1in}

% Custom Command Definitions
\newcommand{\Sball}{\mathcal{S}}
\providecommand{\norm}[1]{\left\lVert#1\right\rVert}

% Hyperref Setup for Clean PDF Links
\hypersetup{
    colorlinks=true,
    linkcolor=blue,
    urlcolor=blue,
    pdftitle={Pattern Arithmetic: Inference from Bowls and Paths},
    pdfauthor={Sandeep Lakshminarayan Chiluveru}
}

% Metadata & Custom Title Banner
\pretitle{\begin{center}\Huge\fontfamily{tt}\selectfont ================================================================================\\[0.5em]}
\posttitle{\\[0.5em] ================================================================================\end{center}}

\title{Pattern Arithmetic: Inference from Bowls and Paths}
\author{Sandeep Lakshminarayan Chiluveru}
\date{\today}

\begin{document}

\maketitle

% ABSTRACT
\begin{abstract}
A companion paper, \emph{Pattern Arithmetic: Typed Algebra of Bounded
Units on a Phase-Colored Sphere}\footnote{Zenodo,
\href{https://doi.org/10.5281/zenodo.22924031}{10.5281/zenodo.22924031},
the concept DOI, which always resolves to the paper's latest
version.}, defines a unit of truth
$x=(r,\hat n,\chi,\tau)$ on a sphere $\mathcal{S}$ and an algebra over
such units --- gather ($\oplus$), remove ($\ominus$), mix ($\boxplus$),
walk ($T$), offset, and assembly ($\bowtie$) --- together with an
admissibility test and an enclosure of external seats. That paper
builds the algebra and stops at a state: a bowl may be mixed, walked,
or assembled and brought to rest. It does not say how a
state, or a bowl, becomes a \emph{decision}.

This paper adds the missing step. A shaman, a tribal meteorologist, or
a market analyst holds a handful of signs and must arrive at one of a
small number of conclusions --- water or no water, storm or calm, rise
or fall. We show that the algebra already contains two distinct
mechanisms for this, corresponding to the two readings the first paper
leaves undemonstrated: an order-free reading, in which evidence
accumulates as phasors and a balanced case returns no answer rather
than a forced one; and an order-dependent reading, in which evidence
is walked as a sequence of actors to a landing compared against
seated names. We give a rule --- the \emph{swap test} --- for which
reading a given body of evidence calls for, and identify four
components a decision needs beyond what the algebra supplies: a
comparison rule, a selection rule for which signs enter and in what
order, a principled abstention, and a place for a lexicon of
conclusions to be learned rather than declared. This paper is a
framework: it specifies these four components precisely enough to be
filled in, and works one example --- the diviner's shells --- by hand.
It does not report an evaluation against data, and says exactly where
that evaluation would have to be made.
\end{abstract}

\section{Introduction}
\label{sec:intro}

A shaman throws ten shells and reads a future from them. He may throw
twenty and read it better, or forty better still --- but unless he
says so, no one else knows what he saw. The shells are evidence; the
saying is the verdict. Pattern Arithmetic gives a precise account of
the shells. It has none, so far, of the saying.

This is deliberate on the first paper's part, not an oversight. That
paper's discipline was to make the algebra trustworthy before asking
it to decide anything: an admissibility test that refuses sorts with
no honest map onto the operators, a type table that says which
combinations are legal, and a running insistence that a walk's landing
is compared to a name, never confused with one. Asking that paper to
also commit to a decision rule would have asked it to defend a claim
it was not ready to make.

It is ready now. This paper takes the atom, the bowl, and the path as
given --- cited, not re-derived --- and asks the next question: given
a bowl of evidence, what is the answer?

\subsection{A human can hold a handful; a machine need not stop
there}
A diviner's working memory bounds how many shells, how many birds, how
many signs can be weighed against one question at once. Nothing in
the algebra of the first paper imposes that bound. The claim of this
paper is narrower than it might sound: not that more evidence is
always better, but that \emph{the mechanism for combining evidence
does not change shape as the amount of evidence grows} --- the same
gather, the same two readings, apply to ten shells or ten thousand
sensor readings. Where that mechanism strains under scale (\S 4.3) is
stated plainly rather than smoothed over.

\section{What Is Inherited}
\label{sec:inherited}

We use, without re-deriving, the following from the companion paper.
Equation and axiom numbers refer to that paper.

\begin{itemize}
\item The atom $x=(r,\hat n,\chi,\tau)\in\mathcal{S}$, with $r\le 1$ a
      normalized strength, $\hat n\in S^2$ a direction, $\chi\in[0,2\pi)$
      a phase, and $\tau$ a sort tag that gates which operators may be
      truthfully applied (Axiom 1; Axiom 0).
\item Gather, $\oplus$: union on finite subsets of $\mathcal{S}$,
      sort-blind, idempotent, lossless, undone by $\ominus$ (Axiom 4).
\item Mix, $\boxplus$: a phasor reading of a bowl whose members share a
      direction and a sort,
      $\boxplus(B)=\bigl(\min(1,|z_B|),\hat n,\arg z_B,\tau\bigr)$
      with $z_B=\sum_{x\in B} r_x e^{i\chi_x}$; commutative,
      non-associative, lossy, uninvertible (Axiom 5).
\item The walk, $T$: an ordered composition
      $T_{w_n}\circ\cdots\circ T_{w_1}(x)$ of actor-parameterized
      rotations and phase shifts; non-commutative in general, and
      many-to-one, so a landing alone does not determine the walk that
      reached it (companion paper, transformation algebra).
\item Assembly $\bowtie$ and the enclosure: two units related while
      both remain intact; external seats $\vec P$ may rest as one
      camp or several. A camp body is measured afterwards; it is not
      a name (companion paper, Axiom~6 and Enclosure).
\item Informational vectors (companion paper, Axiom~7): an atom may
      carry $N$ capped arrows and one uncapped derived sum
      $V_{\mathrm{derived}}$. Coarse assembly and a single-vector
      comparison use that derived pair. The seven diagnostic
      scalars of Axiom~7 are not restated here.
\item \emph{Levels are closed}: a unit carries no representation of
      what it constitutes, and a bowl carries no representation of the
      reading that will be applied to it. Grouping criteria such as a
      time window live outside the 4-tuple.
\end{itemize}

Two consequences of level closure matter enough here to restate. First,
a bowl that saturates under Mix is often a level that has been
flattened --- repeated samples of one unit, mistaken for many units at
one level (companion paper, \S2.7). Second, the atom stays a
four-tuple; grouping criteria such as a time window or a spatial
radius live outside it, deciding bowl membership before the algebra
runs rather than riding inside any single unit.

\subsection{A substitution: bounding aggregated strength}
\label{subsec:bijection_substitution}

The companion paper's Axiom 1 names its clip, $r'=\min(1,|z|)$, as a
choice, and names an order-preserving bijection as the alternative a
setting that aggregates at scale should weigh --- naming this paper's
own use as exactly such a setting, rather than leaving that
identification for the reader to make. This paper takes that
invitation: wherever a strength is reported from a phasor sum for
reading purposes (Reading A's own verdict, and a reference atom's own
reported strength), it substitutes
\begin{equation}
    r' \;=\; \frac{|z|}{K+|z|},
    \label{eq:bijection}
\end{equation}
the $K$-calibrated case of the general Hill form $x^n/(K^n+x^n)$ at
$n=1$, for the companion paper's clip. $K$ is chosen per domain, not
fixed here: it sets where the half-confidence point sits on the scale
of that domain's own evidence. Held at $n=1$, this does not trade away
the reason a bijection was wanted in the first place --- the tail
behaviour $1-r'\approx K/|z|$ for large $|z|$ still lets arbitrarily
large aggregates be distinguished, only at a correspondingly large
scale set by $K$. Raising $n$ above $1$ is not free in the same way:
it steepens the transition near $K$ at the cost of exactly that tail,
reintroducing practical indistinguishability at a smaller scale rather
than removing it, and is not adopted here.

This substitution leaves $\bar R_c$ (Equation~\eqref{eq:coherence})
untouched: that quantity already normalizes by the raw, uncapped
$|z|$, and was built for the same reason this substitution is --- to
stop a bound from erasing a real difference in agreement.

\section{Two Readings of One Bowl}
\label{sec:two_readings}

Given a bowl $B=\{x_1,\ldots,x_n\}$ of evidence bearing on one
question, the companion paper currently defines two readings that turn
$B$ into a single state: phasor Mix and ordered transformation. This is
not a claim that no other reading could be defined; it is the pair
that paper supplies today, and the pair this paper uses. We name them
here because the first paper defines both and demonstrates only one.

\subsection{Reading B: the walk}
\label{subsec:reading_b}
Seat the question as a state $q$. Apply the members of $B$ as actors,
in some order $\pi$:
\[
    \mathrm{land}_\pi(q) \;=\; T_{x_{\pi(n)}}\circ\cdots\circ
    T_{x_{\pi(1)}}(q).
\]
The verdict is the landing, compared against seated names (the storm's
$\mathrm{calm}$ and $\mathrm{storm}$ pins are an instance of this).
Order matters by construction: $T_a T_b\neq T_b T_a$ in general, so
$\mathrm{land}_\pi \neq \mathrm{land}_{\pi'}$ for $\pi\neq\pi'$ is the
expected case, not a failure. This reading is fully worked in the
companion paper's storm example, which we do not repeat.

\subsection{Reading A: the phasor sum}
\label{subsec:reading_a}
Rotate every member of $B$ onto the question's axis $\hat n_q$ (a
rotation, which preserves $r$ and $\chi$), then read the whole bowl at
once:
\[
    \boxplus(B) \;=\; \Bigl(\min(1,|z_B|),\ \hat n_q,\ \arg z_B,\ \tau\Bigr),
    \qquad
    z_B=\sum_{x\in B} r_x e^{i\chi_x}.
\]
This is the companion paper's own clip, stated here as the teaching
cap; \S\ref{subsec:bijection_substitution} substitutes the bijection
wherever this reading is the one actually reported, the shells of
\S\ref{sec:worked_example} among them.
Order plays no role: $z_B$ is a sum, and sums do not care about the
order of their terms. The verdict's \emph{confidence} is $|z_B|$'s
image under the cap, and its \emph{direction} is $\arg z_B$, read
against a lexicon of conclusions seated at particular phases (e.g.\
$\chi=0$ for ``water,'' $\chi=\pi$ for ``no water''). Agreeing evidence
lengthens the sum; opposing evidence shortens it; evenly opposed
evidence returns $z_B=0$ --- black, and a principled abstention rather
than a forced choice (companion paper, Axiom 2).

\subsection{The swap test}
\label{subsec:swap_test}
Which reading applies to a given $B$ is decided by one question: does
exchanging the order in which two members of $B$ arrived change the
answer?
\begin{itemize}
\item If \textbf{no} --- the shells thrown together, ten readings
      taken at one instant --- $B$ has no order to respect, and
      Reading A is the correct mechanism. Applying Reading B here
      manufactures an order-dependence that is not in the evidence.
\item If \textbf{yes} --- birds fleeing, then a wind shift, then a
      temperature drop, where the same three signs in a different
      sequence mean something else --- $B$'s order is itself
      information, and Reading A would discard it. Reading B is
      required.
\end{itemize}
The two readings are not competing hypotheses about one mechanism;
they are the correct tool for two different shapes of evidence, and
the swap test is how a system --- human or otherwise --- tells which
shape it is holding.

\subsection{The readings compose}
\label{subsec:readings_compose}
Nothing prevents using both on one episode. A morning's observations
may arrive in a meaningful sequence (Reading B over the morning) and
also include several simultaneous signs at one moment within it
(Reading A over that moment's bowl, itself one actor in the larger
walk). This is level closure in practice (\S\ref{sec:inherited}): a
Reading A verdict is a state, and a state may act as $w$ in a $T$-path,
exactly as any other seated unit may.

\section{Four Components the Algebra Does Not Supply}
\label{sec:four_gaps}

Both readings above take a bowl and return a state. Neither, on its
own, returns a decision a person or a system can act on. Four
components stand between a state and a verdict, and none of them is
supplied by $\oplus$, $\boxplus$, $T$, or offset.

\subsection{A comparison rule}
\label{subsec:comparison_rule}
The companion paper's storm section says a ball falls \emph{near} the
$\mathrm{calm}$ pin or the $\mathrm{storm}$ pin, and does not define
near. A working comparison rule must state, at minimum: which
coordinates of the landing are read, what metric is used on those
coordinates, and whether $\tau$ must match exactly or may differ.

Two states are compared only if $\tau_A=\tau_B$; $\tau$ gates
comparison the same way it gates every other operator. Given that, a
distance is taken over all three geometric coordinates, each measured
the way the companion paper already measures it elsewhere rather than
by naive coordinate subtraction:
\begin{equation}
    d(A,B)^2 \;=\;
    w_r\,(r_A-r_B)^2
    \;+\;
    w_{\hat n}\,\arccos(\hat n_A\cdot\hat n_B)^2
    \;+\;
    w_\chi\,\bigl(\min(|\chi_A-\chi_B|,\,2\pi-|\chi_A-\chi_B|)\bigr)^2 ,
    \label{eq:comparison_metric}
\end{equation}
using the geodesic distance the offset operator already relies on for
$\hat n$, and the circular distance a phase requires for $\chi$; $r$
alone is a bounded line with no wraparound and needs neither
correction. Raw coordinate subtraction on $\hat n$ or $\chi$ is
\emph{not} a safe default: two directions or phases near $0$ and near
$2\pi$ are adjacent on the sphere or the circle but appear maximally
distant under naive subtraction, and reference atoms placed near a
natural zero point --- which is common, not rare --- would be
misjudged systematically rather than occasionally. \emph{Open in this
paper:} the weights $w_r,w_{\hat n},w_\chi$, meaning how much a
mismatch in strength, direction, or phase should count against an
otherwise close match. \S\ref{sec:worked_example} does not need this
choice made in general: every shell and every reference point there
already shares $\hat n_q$ by construction, so $w_{\hat n}$ never
enters, and the verdict is read from $\chi$ (which reference point is
nearest) with $r$ supplying the null rather than a weighted distance.
A question with reference atoms at different directions would need
$w_{\hat n}$ and $w_\chi$ actually chosen against each other, which
remains open.

\paragraph{A similarity metric.}
$d(A,B)$ answers how far apart two atoms are. The same question
read the other way --- how alike --- is already sitting inside
Equation~\eqref{eq:comparison_metric}. Direction similarity is the
dot product the geodesic term computes before $\arccos$ turns it
into an angle:
\begin{equation}
    \mathrm{sim}(A,B) \;=\; \hat n_A\cdot\hat n_B.
    \label{eq:direction_similarity}
\end{equation}
It is bounded in $[-1,1]$ by construction: $+1$ identical
direction, $0$ orthogonal, $-1$ antipodal. Then
$\arccos(\mathrm{sim}(A,B))$ recovers the geodesic term exactly.
This is the narrow, direction-only reading, the same use cosine
similarity has elsewhere (compare directions, not a
magnitude-weighted blend). It needs none of
$w_r,w_{\hat n},w_\chi$, so it is available today while those
weights stay open. A fuller likeness such as $e^{-d(A,B)}$ would
fold in strength and phase and would inherit the same open
weights; it is not adopted here.

\subsection{A selection rule}
\label{subsec:selection_rule}
For Reading B, the order $\pi$ that the evidence is applied in is
handed to the walk; it is not the walk's to choose. A real system must
decide which signs enter $B$ at all (\S\ref{subsec:comparison_rule}'s
window and radius criteria, companion paper \S2.2) and, when order
matters, in what sequence. This selection is itself where most of an
inference system's actual difficulty lives, and it is entirely outside
the algebra. \emph{Open in this paper.}

\paragraph{A sequence is a degenerate graph.} A walked path
$T_{w_n}\circ\cdots\circ T_{w_1}(x)$ is a simple directed path, one
actor after another with no branching. When $\pi$ is not given in
advance, the natural structure is not a single path but a tree of
candidate paths: several possible next actors branch from a shared
prefix (twinkle, twinkle, then \emph{either} little or something
else), and choosing a walk means searching that tree rather than
replaying one. Every edge of the tree is still one application of
$T$; the tree organizes a search over sequences, it does not add an
operator. How candidates branch, how a branch is scored before it is
fully walked, and how the tree is pruned are the actual content of
this open selection rule, not resolved by naming the structure.

\subsection{A principled null}
\label{subsec:the_null}
Reading A supplies one for free: $z_B=0$. Reading B does not. A walk
always lands somewhere, and a comparison rule (\S\ref{subsec:comparison_rule})
always returns a nearest pin, however far away it is. A Reading-B
system therefore needs an explicit distance threshold beyond which the
verdict is \emph{abstain} rather than \emph{nearest pin regardless}.
This is a second place, distinct from \S\ref{subsec:comparison_rule},
where a number must be chosen rather than derived.

\subsection{Learned pins and learned signs}
\label{subsec:learned_pins}
The companion paper notes that a lookup $E_L$ may be learned
(\S5.6) without the units it points into being anything other than
fixed. The same discipline applies to inference: what is learned is
narrow. A sign's $(r,\chi)$ for a given question --- how strongly, and
in which direction, a particular bird's flight bears on ``storm'' ---
is an $E_L$ entry, learnable from labelled cases. The seated pins
$\mathrm{storm}$, $\mathrm{calm}$ are likewise $E_L$ entries, and may
be learned rather than declared. Neither requires learning a
representation in the sense of an embedding; both are learning one
small map, per sign, per question. \emph{This is the paper's one
concrete claim about the learning loop}: it is cheaper than
representation learning because it is narrow, not because it is
avoided.

\paragraph{Building a reference atom.}
\label{par:reference_atom}
A pin need not be declared; it can be built from witnessed cases the
same way live evidence is read. For a class $c$ (one storm, one
drought), let $G_c=\{y_1,\ldots,y_m\}$ be the Reading-A composite of
each witnessed episode, seated by $E_L$ from that episode's own bowl
and already rotated onto the question's axis $\hat n_q$
(\S\ref{subsec:reading_a}). Because every $y_i$ shares $\hat
n_q$ by construction, $\boxplus(G_c)$ is defined by the
companion paper's bowl-Mix equation ($\S2.7$) without needing its
still-open case of Mix across differing directions (its Future Work);
building a reference atom this way
never depends on a question that paper leaves unresolved. The
reference atom for $c$ is $\mathrm{ref}_c=\boxplus(G_c)$, its strength
reported by Equation~\eqref{eq:bijection} rather than the companion
paper's clip, per \S\ref{subsec:bijection_substitution}.

\paragraph{Coherence, and when to split.}
The capped strength $r(\mathrm{ref}_c)$ is not a safe coherence
signal on its own: it saturates at $\min(1,|z_{G_c}|)$, so a tightly
clustered class and a widely scattered one can both report the
maximal value once enough witnesses push $|z_{G_c}|$ past $1$, at
exactly the point where the underlying spread most needs
distinguishing (\S\ref{par:reference_atom_illustration} makes this
concrete). The uncapped magnitude alone does not fix this either,
since it says nothing about how much evidence produced it: $|z_{G_c}|
= 0.6$ means one thing from a single witness and another from nine
that largely cancelled.

The correct coherence signal normalizes by both at once --- the
\emph{mean resultant length}, the standard measure of angular
concentration for a set of phasors:
\begin{equation}
    \bar R_c \;=\; \frac{|z_{G_c}|}{\sum_{y\in G_c} r_y},
    \qquad
    z_{G_c} \;=\; \sum_{y\in G_c} r_y e^{i\chi_y},
    \label{eq:coherence}
\end{equation}
using the same uncapped $z_{G_c}$ the companion paper's bowl-Mix
equation already computes before capping, alongside one further sum,
$\sum r_y$, tracked once when $G_c$ is assembled and not otherwise
needed for reading live evidence. $\bar R_c\in[0,1]$: $1$ for
perfectly aligned witnesses, $0$ for complete cancellation or uniform
scatter, and --- unlike $r(\mathrm{ref}_c)$ --- unaffected by how many
witnesses there are or how strong each one is, so a class of nine and
a class of ninety are judged by the same scale. This is not a new
invention; $1-\bar R_c$ is the standard circular variance, and the
same quantity underlies tests (Rayleigh's, for one) of whether a set
of angles is distinguishable from uniform scatter at all --- a route
to choosing $\theta$ by a stated significance level rather than by
feel, not adopted here but available.

\paragraph{The direction is a pointer, not only the magnitude.}
\label{par:direction_pointer}
$\bar R_c$ reads the magnitude of $z_{G_c}$; the direction,
$\chi(\mathrm{ref}_c)=\arg(z_{G_c})$, is a separate claim, and it
degrades gracefully rather than failing outright as witnesses
scatter. Simulated against a fixed true direction with $N=8$
witnesses and increasing angular noise, mean error grows smoothly
--- $2.8^\circ$ at low scatter, $15.1^\circ$, $29.1^\circ$,
$50.7^\circ$ at severe scatter --- never a cliff, and continues
rising toward, not past, $90^\circ$, the expected angle to a
genuinely random direction, as scatter grows without bound: the mean
estimate never becomes actively anti-correlated with the truth,
though any single noisy draw still can be. It also sharpens with
more witnesses at fixed noise: $N=2$ gives $36.6^\circ$ mean error,
$N=128$ gives $5.1^\circ$, improving smoothly the way an ordinary
law-of-large-numbers estimator should. None of this is new either
--- $\chi(\mathrm{ref}_c)$ is the circular mean direction, the
direction-side counterpart to $\bar R_c$'s magnitude, both read off
the same $z_{G_c}$.

The two claims must not be conflated. A low $\bar R_c$ does not make
$\chi(\mathrm{ref}_c)$ worthless, only uncertain --- a badly
scattered group can still point roughly toward the truth, the way
$50.7^\circ$ of mean error is still far closer to the truth than to
$90^\circ$ of no information. But a single group's plausible-looking
direction is not, by itself, evidence enough to trust: the same
noise that produces a small average error also produces individual
draws that land unusually close to the truth by chance,
indistinguishable, from one reading alone, from a genuinely coherent
source. What separates the two is repetition --- several independent
groups' directions agreeing with each other is exactly the claim
$\bar R_c$ already tests, applied one level up, over a bowl of
directions rather than a bowl of witnesses.

Given a coherence threshold $\theta\in(0,1)$ and a minimum
group size $k_{\min}$ below which a split is not attempted: if
$\bar R_c\ge\theta$, keep $\mathrm{ref}_c=\boxplus(G_c)$; otherwise, if
$|G_c|\ge k_{\min}$, partition $G_c$ into sub-groups and recurse on
each, so that a single label may resolve into several reference atoms
(a slow-onset and a flash drought, say); otherwise accept
$\mathrm{ref}_c$ as weak, or fall back to comparing incoming evidence
against $G_c$'s individual members rather than their compression. How
$G_c$ is partitioned when a split is warranted is not specified here
and is left open, as is the choice of $\theta$ itself; $\theta$ must
be bounded away from $0$, since every singleton group has $\bar
R_c=1$ trivially and a threshold with no floor never stops splitting.
Group size strictly decreases at each split, so the recursion
terminates regardless of $\theta$'s exact value.

\paragraph{This is level closure, applied to references.} A class
discovered to need splitting is a level-2 object in its own right
(companion paper, \S2.2): the split is discovered from $G_c$'s own
internal agreement, not asserted in advance, and each resulting
sub-reference is exactly as usable --- comparable, walkable, further
splittable --- as the class it came from. A domain with genuinely
few witnesses (two droughts, not nine) is not a failure of the
procedure; $k_{\min}$ says plainly that there is not yet evidence to
subdivide responsibly, and the honest output is a weaker reference or
a fall back to individual comparison, not a manufactured split.

\paragraph{Repeated splitting is a tree, and classifying against it
means descending it.} A label that splits once and then splits again
within a child produces a tree of reference atoms rooted at the
original class: $\mathrm{storm}$ at the root, perhaps $\mathrm{slow\
onset}$ and $\mathrm{flash}$ as children, each itself subject to the
same $\bar R\ge\theta$ test on its own witnesses. Comparing an
incoming atom against such a tree is naturally top-down: match against
the root's siblings first (storm, drought, rain, flood), then, within
whichever is nearest, descend to its children if it has any, and
repeat. This is the shape of a classical decision tree, but not its
mechanism: a decision tree tests a raw coordinate against a learned
threshold at each node; a node here compares a full state against a
reference atom by Equation~\eqref{eq:comparison_metric}, using
operators and a metric the algebra already supplies rather than
introducing threshold tests as a new primitive. The tree is discovered
by $\bar R$, level by level, the same way \S\ref{par:reference_atom}
discovers it, not designed in advance; the comparison rule and the
null of \S\ref{subsec:comparison_rule} and \S\ref{subsec:the_null}
apply at every node unchanged.

\paragraph{Merging groups.} Two witnessed groups of the same $\tau$,
of any sizes, merge exactly, without revisiting either group's raw
witnesses: $z_{G_c}$ and $\sum r_y$ are sufficient statistics, in the
standard sense --- each additive over a disjoint union, so
\begin{equation}
    z_{\text{merged}} = z_{G_1} + z_{G_2},
    \qquad
    \Bigl(\sum r_y\Bigr)_{\text{merged}}
    = \Bigl(\sum r_y\Bigr)_{G_1} + \Bigl(\sum r_y\Bigr)_{G_2},
    \label{eq:sufficient_stat_merge}
\end{equation}
give $\bar R_{c,\text{merged}}$ and $\mathrm{ref}_{\text{merged}}$
identical to recomputing Equation~\eqref{eq:coherence} from the two
groups' combined raw membership --- checked to full floating-point
precision on groups of size $5$ and $9$. Group size needs no special
handling: it is already carried inside $\sum r_y$, so a group of
$400$ correctly dominates a group of $3$ in the merged result without
any explicit weighting step. The cost of merging is therefore $O(1)$
in the number of witnesses, holding only the two summaries rather
than either group's membership.

Two witnessed groups of \emph{different} $\tau$ do not merge this
way --- averaging a storm phasor sum against a drought one asserts a
relationship comparison and Mix both already refuse, since both are
$\tau$-gated. The companion paper's own Axiom 4 already supplies the
right mechanism for this case: gather is sort-blind, so the two
groups combine losslessly into one bowl regardless of sort,
preserving both types distinctly rather than fusing them; any further
reading --- $\bar R_c$, $\mathrm{ref}_c$, Equation~\eqref
{eq:sufficient_stat_merge} --- is then applied separately within each
$\tau$ present in the combined bowl, exactly as it already is for any
mixed-sort bowl. Same-sort merging needs a new identity because $\bar
R_c$ is new to this paper; cross-sort merging needed none, since
$\oplus$ was never sort-restricted to begin with.

\paragraph{Illustration.}
\label{par:reference_atom_illustration}
Nine witnessed episodes at $r=0.6$ each, at
$\chi=8^\circ,9^\circ,9^\circ,10^\circ,10^\circ,10^\circ,11^\circ,
11^\circ,12^\circ$, give $\sum r_y=5.4$ and
$\mathrm{ref}=(1.00,\hat n_q,10.0^\circ)$: capped $r'=1.00$ and
$\bar R=0.9998$ agree --- coherent, confidently kept. The same nine
witnesses spread instead across
$\chi=-40^\circ,-20^\circ,-10^\circ,0^\circ,10^\circ,20^\circ,
30^\circ,45^\circ,60^\circ$ give the identical capped $r'=1.00$, since
$|z_{G_c}|$ still exceeds the cap --- but $\bar R=0.8689$, correctly
showing this class as less coherent than the first even though the
capped reading cannot tell them apart. Five witnesses at $\chi=0^\circ$
and four at $\chi=180^\circ$ --- two real sub-types wearing one
label --- give capped $r'=0.60$ and $\bar R=0.11$: both signal a
problem, but $\bar R$ states it as a small fraction of the evidence
actually agreeing, not merely a moderate reading. A perfectly even
split, five and five, gives $\bar R=0.00$, the same cancellation that
gives \S\ref{sec:worked_example} its abstain, now flagging the class
itself rather than a single reading of it.

\subsection{A walk-based classifier}
\label{subsec:walk_classifier}

Everything above classifies by comparing a state against a
\emph{compressed} reference built via $\boxplus$. The same question
--- what does new evidence resemble --- has a second mechanism,
using Reading B instead of Reading A, that has been available since
\S\ref{subsec:reading_b} and simply never assembled for this purpose:
walk new evidence to a landing, $x=\mathrm{land}_\pi(q)$, and compare
that landing against a stored reference by
Equation~\eqref{eq:comparison_metric} --- exactly what the companion
paper's storm section already gestures at (``an endpoint may fall
near the storm pin or the calm pin''), now made precise by the metric
\S\ref{subsec:comparison_rule} supplies rather than left as an
undefined ``near.'' Nothing here is new machinery; it reuses Reading
B and the comparison rule exactly as they already stand.

\paragraph{Why this classifier does not compress.}
What it cannot do, for now, is combine several witnessed landings
into one reference the way \S\ref{par:reference_atom} combines
several Reading-A composites into $\mathrm{ref}_c$. Every $y_i$
feeding a Reading-A reference atom shares $\hat n_q$ by construction
--- Reading A's own rotation step guarantees it
(\S\ref{subsec:reading_a}). A $T$-walk's landing carries no such
guarantee: nine witnessed episodes, walked independently through
their own sequences of actors, generally land nine different places.
Compressing those landings via $\boxplus$ would need exactly the
companion paper's still-open case of Mix across differing directions
(its Future Work), or a combination rule for rotations this algebra
does not currently define --- averaging rotations is a real, known
operation (a Karcher mean on $\mathrm{SO}(3)$, say) but importing one
is a decision for later, not assumed here.

\paragraph{The fallback is the whole mechanism, not a special case.}
A walk-based classifier therefore keeps every witnessed landing
individually and compares a new landing against each --- the same
fallback \S\ref{par:reference_atom} already names for a class too
small to split (``fall back to comparing incoming evidence against
$G_c$'s individual members''); here it is the only available
mechanism rather than a fallback from a failed compression. A single
declared pin, as in the storm example, is simply the case of one
witness kept rather than several.

\paragraph{Illustration.} Seat the question at the north pole,
$\hat n_q=(0,0,1)$. Two actors walk it: $r_{w_1}=\tfrac13$
($\alpha=60^\circ$) about the $x$-axis, then $r_{w_2}=\tfrac13$
($\alpha=60^\circ$) about the $y$-axis, landing at
$\hat n=(0.433,-0.866,0.25)$, i.e.\ $(\theta,\phi)=(75.5^\circ,
-63.4^\circ)$. Two declared pins: $\mathrm{dry}$ at
$(\theta,\phi)=(80^\circ,300^\circ)$ and $\mathrm{rain}$ at
$(\theta,\phi)=(30^\circ,90^\circ)$. The geodesic distance
(Equation~\eqref{eq:comparison_metric}, $w_r=w_\chi=0$) from the
landing to $\mathrm{dry}$ is $5.6^\circ$; to $\mathrm{rain}$,
$102.5^\circ$. The verdict is $\mathrm{dry}$, by a wide and
unambiguous margin, using nothing beyond a walk already defined and a
metric already stated.

\emph{Blocked for later, not solved here:} a genuine compressed
reference path, standing in for several witnessed walks the way
$\mathrm{ref}_c$ stands in for several Reading-A composites, needs
either the cross-direction Mix question resolved or an explicit
rotation-averaging operator adopted and justified on its own terms.
Both are left for iteration.

\section{Worked Example: The Diviner's Shells}
\label{sec:worked_example}

We work Reading A in full, since the companion paper never carries a
phasor reading past two states, and its property demonstration
(non-associativity, on an anonymous triple) is not a reading of a
bowl that means something.

Seat the question \emph{will it rain} at $\hat n_q$. Each shell, once
cast, is read by a diviner (or, later, by a learned $E_L$) as a state
$(r,\chi)$ at $\hat n_q$: $\chi=0$ favors rain, $\chi=\pi$ favors no
rain, and intermediate $\chi$ is a hedge. Five casts, each shell
carrying $r=0.4$ unless noted:

\begin{center}
\begin{tabular}{lcc}
\hline
Cast & reading & interpretation \\
\hline
4 shells, all favor rain      & $r=1.00$ at $\chi=0^\circ$   & confident rain \\
3 favor, 1 against            & $r=0.80$ at $\chi=0^\circ$   & rain, less confident \\
2 favor, 2 against            & $r=0.00$                     & \textbf{abstain} \\
4 favor, 2 hedge at $90^\circ$ & $r=1.00$ at $\chi\approx 26.6^\circ$ & rain, hedged direction \\
12 shells ($r=0.15$ each), all weakly favor    & $r=1.00$ at $\chi=0^\circ$   & confident rain \\
\hline
\end{tabular}
\end{center}

The third row is the case worth dwelling on. Two shells favor rain and
two oppose it, at equal strength; the phasors cancel exactly, and the
verdict is black --- not a coin flip, not a default to the majority
class, but the algebra's own statement that this evidence does not
decide the question. A system built on a forced classifier (softmax
over two classes, say) has no way to express this outcome; it must
output a confident-looking number even when the evidence is silent.
Reading A's null is not an afterthought fitted onto the algebra; it is
Axiom 2, doing here exactly what it does for a Mix reading of colour.

The fifth row shows the saturation the companion paper's \S2.7
discusses: twelve weak, agreeing signs pin at $r=1$ exactly as four
strong ones do, and the difference between ``moderately confident''
and ``overwhelmingly confident'' is lost. This is the place where
\S\ref{subsec:selection_rule}'s selection rule matters most: whether
those twelve readings are twelve units at one level, or repeated
samples of one unit that should have been reduced (companion paper,
Axiom 4) before ever reaching the bowl, is a question this paper does
not resolve and flags as the clearest candidate for the next.

\S\ref{subsec:bijection_substitution}'s substitution softens this
specific symptom without resolving the question just raised.
Calibrated to this domain's evidence with $K=0.2$, the fourth row
reads $r'=0.899$ and the fifth $r'=0.900$ --- genuinely
distinguishable, where the clip read both as $1.00$ --- and even a
hundred thousand such weak witnesses would read $r'=0.99999$, not
exactly $1$. What is now settled is that the two \emph{can} be told
apart; what remains open is whether twelve weak signs are honestly
twelve co-present witnesses of one thing, or repeated samples of a
single witness that should have been reduced before reaching the bowl
at all. The substitution reports strength faithfully either way and
does not decide which reading was owed.

\emph{Reading B on the same shells} is not worked here. Doing so
honestly requires an order for the casts, which a physical handful of
shells thrown at once does not supply --- exactly the case the swap
test (\S\ref{subsec:swap_test}) assigns to Reading A. A worked Reading
B example belongs to evidence that genuinely arrives in sequence, and
the companion paper's storm already provides one. A smaller one,
native to this paper, is still worth having: three actors show the
same point without the storm's six steps.

Seat two actors, $\mathrm{twinkle}$ with axis $\hat n=(1,0,0)$ and
strength $r=0.5$ (a quarter turn, $\alpha=90^\circ$), and
$\mathrm{little}$ with axis $\hat n=(0,1,0)$, the same strength.
Start the occupant at the north pole, $x_0=(0,0,1)$. Applying
$\mathrm{twinkle}$ first sends it to $(0,-1,0)$, a point that now
lies exactly on $\mathrm{little}$'s own axis, so $\mathrm{little}$
turns it nowhere further: $\mathrm{twinkle}$ then $\mathrm{little}$
lands at $(0,-1,0)$. Applying $\mathrm{little}$ first sends $x_0$ to
$(1,0,0)$, exactly on $\mathrm{twinkle}$'s own axis this time:
$\mathrm{little}$ then $\mathrm{twinkle}$ lands at $(1,0,0)$. Two
different pins, $(\theta,\phi)=(90^\circ,270^\circ)$ against
$(90^\circ,0^\circ)$, from the same two actors --- order the only
thing that changed. This is the companion paper's own
non-commutativity (\S3.2) made concrete at three steps rather than
argued in the abstract, and each half of it is a single rotation a
reader can check without software: the second step, in both orders,
is a vector rotating about itself.

\section{Level Closure in Practice: Correcting a Path}
\label{sec:level_closure_practice}

A $T$-path records the sequence of actors applied to it, faithfully,
by definition (companion paper, Axiom 6). It has no means of knowing
whether that sequence was the right one --- level closure
(\S\ref{sec:inherited}) says precisely this: a lower unit carries no
representation of what it constitutes, and neither the atom nor the
path can see the higher-level context that would tell it it has erred.

A concrete case: a sequence of observed signs is walked in the order
observed, but a later sign (or a completed higher-level unit, such as
the rest of a sentence or the rest of a song) makes that order appear
wrong --- \emph{twinkle, little, twinkle, star} rather than
\emph{twinkle, twinkle, little, star}. Correcting this is a level-2
operation: it requires comparing the recorded path against context the
path itself does not have access to, and revising the walk
probabilistically. This paper's position is that the path must
therefore be \emph{legible} to the level above --- readable as an
ordered record, not only walkable --- so that a corrector at the next
level can inspect and revise it. Nothing in the companion paper's
axioms prevents this; nothing in them provides it either. We flag it
as the clearest point of contact between this paper's verdicts and a
learning loop operating above them, and leave the corrector itself
unspecified.

\section{Prediction}
\label{sec:prediction}

Everything above compares an atom to a reference built from atoms of
the \emph{same} concurrent identity: a $\mathrm{storm}$ reference is
built from patterns that were themselves storms, and a match says what
the evidence currently resembles. That is classification. Prediction
--- what is likely to happen, not what this already is --- needs a
different labelling convention, not a different mechanism.

\subsection{Immediate forecast: relabelling by outcome}
\label{subsec:immediate_forecast}

Fix a lead time $\Delta$. A \emph{precursor witness} for outcome $c$ is
a Reading-A composite $y_i$ of the evidence bowl observed at some time
$t$, labelled not by what it is but by what followed it: $c$ occurred
at $t+\Delta$. Gathering such witnesses gives a precursor group
$G_c^{\Delta}$, and everything in \S\ref{par:reference_atom} applies
unchanged: $\mathrm{ref}_c^{\Delta}=\boxplus(G_c^{\Delta})$, checked for
coherence by $\bar R_c$ (Equation~\eqref{eq:coherence}) and split if
$c$ conceals more than one precursor shape, exactly as
\S\ref{par:reference_atom_illustration} splits $\mathrm{drought}$.
Forecasting is then comparison against $\{\mathrm{ref}_c^{\Delta}\}$ by
the ordinary rule of Equation~\eqref{eq:comparison_metric}: the
nearest precursor reference, within threshold, says what is likely at
$t+\Delta$. The null of \S\ref{subsec:the_null} applies exactly as
before --- no close precursor match is an honest \emph{no forecast},
not a forced one, which matters most exactly where a wrong call is
costly.

Nothing here is a new equation. The only change from
\S\ref{subsec:learned_pins} is which witnesses are gathered under a
label: concurrent identity for classification, a fixed time-shifted
outcome for this. Two things mark it as a first, crude pass rather
than a finished one, named so they are not mistaken for settled
choices: $\Delta$ is fixed and single-valued, so the forecast has one
horizon and no notion of a nearer or more distant event both being
plausible; and no forecast carries a horizon-dependent confidence ---
a match found for $\Delta=10$ minutes and one for $\Delta=10$ days are
read with the same $\bar R_c$ scale, though the latter should
plausibly be trusted less for the same coherence value. Both are left
for iteration, not resolved here.

\subsection{A cheap decoder: transition counting over actors}
\label{subsec:markov_decoder}

Everything above answers ``what does this resemble.'' A different
question --- ``what comes next in an unfolding walk'' --- has a
mechanism simpler than either technique in
\S\ref{subsec:four_cells}, and it is worth stating precisely because
it needs none of the machinery either one does: no metric, no
$\boxplus$, no coherence check. It is a first-order Markov chain over
the discrete vocabulary $E_L$ already provides --- the same technique,
under the same name, used since the earliest work on statistical
language modelling.

Given a set of witnessed trajectories, each already an ordered
sequence of actors $(w_1,\ldots,w_n)$ (companion paper's
pipeline equation, \S2.2), count how often each actor is followed
by each other across every witnessed path:
\begin{equation}
    \widehat{P}(w_{k+1}=w' \mid w_k=w)
    \;=\;
    \frac{\mathrm{count}(w\to w' \text{ across all witnessed paths})}
         {\sum_{w''}\mathrm{count}(w\to w'' \text{ across all witnessed paths})}.
    \label{eq:markov}
\end{equation}
Decoding is a lookup: given the actor just visited, predict the $w'$
maximizing $\widehat P$, or sample from it. Nothing here touches
$T_w(x)$; the mechanism operates entirely on \emph{which} actor
follows \emph{which}, never on the rotation either one performs. It
is closer to bookkeeping on the tape a walk already produces than to
a use of the walk's own geometry.

\paragraph{Illustration.} Four witnessed paths ---
$(\text{twinkle},\text{twinkle},\text{little},\text{star})$ twice,
$(\text{twinkle},\text{little},\text{star})$ once, and
$(\text{twinkle},\text{twinkle},\text{little},\text{bat})$ once --- give
\[
\widehat P(\cdot\mid\text{twinkle}) =
\Bigl\{\text{twinkle}:\tfrac{3}{7},\ \text{little}:\tfrac{4}{7}\Bigr\},
\qquad
\widehat P(\cdot\mid\text{little}) =
\Bigl\{\text{star}:\tfrac{3}{4},\ \text{bat}:\tfrac{1}{4}\Bigr\}.
\]
A single witnessed path alone would have given an uninformative
$\tfrac{1}{1}$ split; the distinction above only exists because
several paths were witnessed, not one.

\paragraph{What this mechanism does not use, stated plainly.} The
atom carries $r$, $\hat n$, $\chi$; this decoder reads none of them,
only $\tau$-gated discrete identity. For this specific mechanism, a
bare token would do exactly as well --- the geometry is present but
idle. It is named here rather than left implicit, since claiming
otherwise for this particular technique would overstate what it
uses.

\paragraph{Cost, stated honestly.} It needs several witnessed
repetitions before its output means anything, and it only ever looks
back a short, fixed window (one prior actor here; an $n$-th order
chain extends the window at the cost of needing correspondingly more
witnessed data to fill it) --- it has no notion of the deeper
structure a full prefix match, \S\ref{subsec:path_predictors} below,
would catch. Between the two, this is the cheaper, unblocked
mechanism for the same job; the exemplar-based one is the deeper,
still-open alternative.

\paragraph{Where this sits against a standard pipeline.} The
companion paper's own stack (Related Work: Transformers) is
\[
\text{sentence}\to\text{token}\to E_L\to\text{sphere}\to
\oplus\text{ (bowl) }+\;T\text{ (walk)}.
\]
This decoder extends that diagram along an axis it does not yet
have. The diagram's own $T$ assumes its actor given --- one
witnessed sequence, walked. This mechanism is what supplies a next
actor when none is given: a policy read from witnessed sequences,
not a rotation computed from one. It does not stand where $T$
stands --- nothing above computes $T_w(x)$, as already stated ---
it stands one step earlier, choosing which $w$ a
walk hands to $T$ next:
\[
\oplus\text{ (bowl)}\ \to\ \text{transition table}\ \to\ T\text{ (walk)},
\]
not a fork between two fillers of one slot. Two distinct claims
follow from adding this step, and they should not be run together.
\emph{Drop-in address}: wherever a downstream pipeline needs one
vector per token, $E_L$ can occupy that slot; an atom's $(r,\hat
n)$, or a listed inner bundle, reads where a learned
high-dimensional embedding otherwise would. \emph{Same verbs,
foreign net}: a transformer may still be what computes $T_w(x)$'s
rotation, or what estimates the transition table feeding it --- two
different jobs, both unreplaced, each doing what it already does on
an address this paper supplies rather than one it learns. The claim
in either case is a shared address and a shared verb set, never that
one technique outperforms the other on some stated task; nothing
here has measured that, and nothing here claims to.

The first use asks more of $E_L$ than the second, and the difference
is worth naming precisely rather than leaving implicit.
\S\ref{subsec:learned_pins}'s own claim is that entries are
\emph{learnable}; it is silent on whether two implementations
learning entries independently are learning the same lexicon at all.
A genuine drop-in needs more than that: $E_L$ must be one map,
published once against a stated vocabulary --- a specific word list
or dictionary edition another system can point at --- not a private
assignment remade per use. Absent that, ``one spelling, one atom''
holds within a single run and not across two, and the drop-in use
quietly collapses into the second: a private notation a transformer
might still fill, not an address a second system could read. This
paper does not supply that agreement; it names the requirement so
the first use is not claimed before it is met.

One admissibility point carries over unchanged, and is worth stating
for this mechanism specifically rather than left to be assumed.
Axiom 0's $\mathtt{num}$ discipline --- decode, add in $\mathbb{Z}$,
encode, never Mix --- governs a sort regardless of which mechanism
above is reading it: a chain walked or counted over mixed
$\mathtt{word}$ and $\mathtt{num}$ actors keeps that refusal, and
does not Mix two numeric tokens merely because both landed in the
same trained table. Being learned rather than declared exempts
neither this decoder nor a corpus-trained $T$ from it.

\subsection{Open: path predictors}
\label{subsec:path_predictors}

The immediate forecast reads only the endpoint of what happened after
a precursor; a deeper forecast reads the whole unfolding. Store, per
witnessed episode, the entire trajectory that followed --- not one
composite but the ordered sequence of atoms, timestamped as the
companion paper's step-metadata generalization allows (\S3.3) --- and
match an
\emph{observed prefix} of a trajectory now against the opening steps of
stored trajectories, reading the continuation of the closest match as
the prediction. This is not a new idea invented for this paper: it is
the analog method already used in operational forecasting, arrived at
here from the algebra's own machinery rather than borrowed from it.

It is filed open, not attempted, because it needs two things this
paper has not built. First, storage of complete historical
trajectories rather than single compressed reference atoms, a real
cost the immediate forecast does not pay. Second, and more
substantial: Equation~\eqref{eq:comparison_metric} compares one state
to one state, and prefix matching compares a \emph{sequence} of states
to a sequence, which needs a distance over sequences --- some
alignment or step-wise accumulation of the single-state metric,
robust to two episodes unfolding at different paces --- that has not
been defined. This section marks where that work belongs rather than
attempting it; the immediate forecast of
\S\ref{subsec:immediate_forecast} does not depend on it and should not
wait for it.

\subsection{Four cells, two techniques}
\label{subsec:four_cells}

Comparison (\S\ref{subsec:comparison_rule}) is the one operation
behind everything in this section and the last --- everything, that
is, except \S\ref{subsec:markov_decoder}'s transition count, which
answers a different question (what comes next, not what this
resembles) without comparison at all. Two established
techniques are built from comparison, depending on what a produced
state is checked against: \emph{nearest-centroid classification},
comparing against a single reference compressed from many witnesses
via $\boxplus$ (\S\ref{par:reference_atom}), and
\emph{nearest-neighbor classification}, comparing against witnesses
kept individually (\S\ref{subsec:walk_classifier}). Neither is
specific to this paper --- both are standard, decades-old techniques
in pattern recognition, arrived at here from the algebra's own
constraints rather than imported by name.

A second, independent axis is what the label means: identity, giving
a verdict about now, or a time-shifted outcome, giving a forecast
about later (\S\ref{subsec:immediate_forecast}). The two axes cross
to give four cells, not a new technique for each:

\begin{center}
\begin{tabular}{lll}
\hline
 & now (identity label) & later (outcome label) \\
\hline
nearest-centroid &
  \S\ref{par:reference_atom} &
  \S\ref{subsec:immediate_forecast} \\
nearest-neighbor &
  \S\ref{subsec:walk_classifier} &
  \S\ref{subsec:path_predictors} \\
\hline
\end{tabular}
\end{center}

The bottom-right cell already has its own established name for the
same reason the top row does: nearest-neighbor comparison, run
against outcome labels, over whole sequences rather than single
states, is the analog method \S\ref{subsec:path_predictors} cites.

The split between rows is not stylistic; it is forced by the same
fault line as $\boxplus$ and $T$. Compression into a centroid is
available wherever $\boxplus$'s composites are guaranteed to share a
direction; keeping individual neighbors is what remains wherever that
guarantee fails --- a $T$-landing (\S\ref{subsec:walk_classifier}),
or, in the one open cell, a whole sequence with no defined distance
yet.

Three of the four cells are specified and checked on worked numbers;
none should be read as more finished than that. Nearest-centroid
classification (\S\ref{par:reference_atom}) is complete as a
mechanism, with three choices left named and open within it: the
coherence threshold $\theta$, the comparison weights
$w_r,w_{\hat n},w_\chi$, and how a class is partitioned once a split
is warranted. Nearest-neighbor classification
(\S\ref{subsec:walk_classifier}) is complete for what it does ---
walking new evidence and comparing the landing against kept
witnesses --- with compression into a single centroid blocked
separately, not left unfinished. The immediate forecast
(\S\ref{subsec:immediate_forecast}) is explicitly a first pass: one
fixed lead time $\Delta$, no horizon-dependent confidence, both named
for iteration rather than resolved. The fourth cell,
\S\ref{subsec:path_predictors}, is open outright, waiting on the
sequence-level distance neither this table nor that section
supplies. None of the three built cells has been evaluated against
real data; each is verified on the toy numbers of its own
illustration, per the framework-versus-benchmarked choice
\S\ref{sec:scope} leaves open.


\section{What this paper does not redo}
\label{sec:not_redo}

The companion paper now states a fifth verb, assembly ($\bowtie$),
and an enclosure in which external seats $\vec P$ may descend an
energy until they rest as one camp or several
(companion paper, Axiom~6 and the Enclosure section). Those
constructions are cited, not repeated. A camp body is not a verdict.
Naming a resting camp still requires the comparison rule of
\S\ref{subsec:comparison_rule} against a seated pin. Rotate-then-Mix
in Reading~A is a convention of this paper: it does not close the
companion's open question of Mix across differing directions. What
assembly itself settles, and what it leaves open, is no longer
tracked here: a second companion paper, \emph{Pattern Arithmetic:
Higher Truths from Assembly}, now carries that question forward on
its own.

\subsection{Multiplicity as a reading, not a rewrite of $\oplus$}
\label{subsec:multiplicity_reading}
The companion paper's $\oplus$ is set union: $A\oplus A=A$. When
inference must keep how many independent witnesses produced the same
tuple, this paper reads a bowl through a multiplicity
$m:\mathcal{S}\to\mathbb{N}_0$ alongside the set, never in place of
it. $m(A)=2$ is two witnesses of one unit, not a second seat and not
a change to Axiom~4. Aggregates that need count --- $\sum r$,
$\bar R_c$ --- use $m$. Presence questions continue to use $\oplus$.

\section{Scope and Open Questions}
\label{sec:scope}

This paper is a framework: mechanism plus one worked toy example
(the shells). It does not report an evaluation against data.

The following remain open rather than resolved here:

\begin{itemize}
\item The comparison weights $w_r,w_{\hat n},w_\chi$
      (\S\ref{subsec:comparison_rule}).
\item The selection rule (\S\ref{subsec:selection_rule}), including
      the companion paper's grouping window.
\item Saturation under aggregation (\S\ref{sec:worked_example}):
      whether to reduce repeated samples before the bowl, or rely on
      the bijection of \S\ref{subsec:bijection_substitution}.
\item The path corrector (\S\ref{sec:level_closure_practice}).
\item Path predictors (\S\ref{subsec:path_predictors}); the immediate
      forecast does not require them.
\item The coherence threshold $\theta$ and the partition rule once a
      class splits (\S\ref{par:reference_atom}).
\item True angular dispersion: distinguishes genuine scatter from a
      tight two-cluster split, which $\bar R_c$ and $z_{G_c}$ both
      read as zero either way. Confirmed as a real, surviving blind
      spot --- tested against $\alpha_{\text{atom}}$ (companion
      paper, Axiom~7) and not closed by it either ($-0.43$ scattered
      versus $-0.50$ bimodal, too near to separate).
\end{itemize}

What this paper delivers is narrower than a solved problem and wider
than a single technique: two readings, a rule for choosing between
them, and four components named precisely enough that filling one in
does not require touching the others. The four are not evenly
finished. A comparison rule and a principled null are complete down
to a stated calibration; a selection rule is characterized rather
than built, and said to be where the real difficulty lives, not
hidden as if it were not. A verdict, once reached, is a state like
any other and may re-enter the algebra as an actor in a further
walk (\S\ref{subsec:readings_compose}) --- the mechanism this paper
adds does not stop at one decision; it hands its own output back to
the units it was built from.

\end{document}
